\documentclass[11pt]{article}
% Standalone preamble for web-only supplementary notes.
\usepackage[T1]{fontenc}
\usepackage{amsmath,amssymb,amsthm}
\usepackage{hyperref}
\usepackage{cleveref}
\DeclareMathOperator{\Map}{Map}
\DeclareMathOperator{\Fun}{Fun}
\newcommand{\iso}{\xrightarrow{\sim}}
\newtheorem{lemma}{Supplementary Lemma}
\newtheorem{corollary}[lemma]{Supplementary Corollary}
\crefname{lemma}{Supplementary Lemma}{Supplementary Lemmas}
\Crefname{lemma}{Supplementary Lemma}{Supplementary Lemmas}
\crefname{corollary}{Supplementary Corollary}{Supplementary Corollaries}
\Crefname{corollary}{Supplementary Corollary}{Supplementary Corollaries}

\begin{document}
\section*{Mapping out of pasted squares}

The \href{pushouts.html\#lem:Pasting_Lemma_Pushouts}{pushout pasting proof}
applies the pullback pasting lemma after mapping out.
This requires comparing the specified matching identifications on the two
routes. The calculation uses the Chapter 1 mapping-anima and pullback
structure; functor categories are not needed.

Consider adjacent squares $L,R$ with top arrows
$g_1\colon A_1\to A_2$, $g_2\colon A_2\to A_3$, bottom arrows
$h_1\colon B_1\to B_2$, $h_2\colon B_2\to B_3$, vertical arrows
$f_i\colon A_i\to B_i$, and matchings
\[
 \lambda\colon f_2g_1=h_1f_1,\qquad
 \rho\colon f_3g_2=h_2f_2.
\]
Their pasted square $O$ uses the matching from $f_3(g_2g_1)$ to
$(h_2h_1)f_1$ obtained by reassociating, applying $\rho$ along $g_1$,
reassociating, applying $h_2$ to $\lambda$, and reassociating back.
The matching is part of the input, as well as the four boundary functors.

\begin{lemma}
\label[lemma]{lem:Supplementary_Product_Pasting}
For an absolute category $X$, put $L_X(f)=\mathrm{id}_X\times f$, and
write $\kappa_X(f,g)\colon L_X(g)L_X(f)=L_X(gf)$ for its chosen composition
comparison. The product images of $L,R$ paste to a cocone which compares
with the product image of $O$, after changing the upper span arrow by
$\kappa_X(g_1,g_2)$. The comparison has the identity on the left leg and
$\kappa_X(h_1,h_2)$ on the right leg, with compatibility of the matchings.
\end{lemma}
\begin{proof}[Proof idea]
Expand the pasted matching into its five stages. Compatibility with inverse
associators handles the first and last stages. Naturality of the product
composition comparison handles the two whiskered matchings, and its
associativity comparison handles the middle stage. Paste these five
comparison squares and cancel the common outer composition comparison.
Postcomposition by any functor $X\times B_3\to E$ preserves the
resulting whole-cocone comparison.
\end{proof}

For $k\colon A\to B$, write $k^*\colon\Map(B,E)\to\Map(A,E)$
for precomposition by $k$.

\begin{lemma}
\label[lemma]{lem:Supplementary_Mapping_Pasting}
For an absolute category $E$, let $M_E(L)$, $M_E(R)$, and $M_E(O)$ be
the mapping-out cones. Pasting $M_E(R)$ and $M_E(L)$, in this reversed
order, and changing its left cospan arrow by
\[
 \kappa_g\colon g_1^*g_2^*=(g_2g_1)^*
\]
produces a cone with a specified comparison to $M_E(O)$.
This comparison includes its matching equation, without any pushout
hypothesis on $L$ or $R$.
\end{lemma}
\begin{proof}[Proof idea]
Use the universal parameter $X=\Map(B_3,E)$ and uncurry the mapping-out
cones at its identity. The composition comparison for precomposition on
mapping animae evaluates to the product comparison $\kappa_X$.
The two square matchings evaluate to the corresponding product cocones,
postcomposed with evaluation $X\times B_3\to E$.

Apply \Cref{lem:Supplementary_Product_Pasting} and reflect the resulting
whole-cocone comparison through uncurrying. This reflection is available
because $X$ is an anima. Finally remove restriction by the identity with
the cone unit comparison.

If $U,V$ denote the source and target cones over a common cospan $(p,q)$,
and $a,b$ the constructed leg comparisons, the retained equation is
\[
 \operatorname{match}(V)\cdot(p\star a)
   =(q\star b)\cdot\operatorname{match}(U).
\]
Here $\star$ denotes whiskering and vertical composition applies the
rightmost factor first. The witnesses $a,b$ are those constructed by
uncurrying and reflection; the assertion includes their compatibility.
\end{proof}

\begin{corollary}
\label[corollary]{cor:Supplementary_Pushout_Pasting}
If $L$ is a pushout, then $R$ is a pushout if and only if the specified
pasted square $O$ is a pushout.
\end{corollary}
\begin{proof}[Proof idea]
For every $E$, the mapping-out cone $M_E(L)$ is a pullback. Contravariance
makes it the fixed right square in the pullback pasting argument.
Pasting and cancellation therefore identify the pullback conditions for
$M_E(R)$ and the pasted mapping cone. Changing the cospan arrow by
$\kappa_g$ preserves the pullback property.
\Cref{lem:Supplementary_Mapping_Pasting} transfers it to $M_E(O)$,
including the specified matching. Quantify over $E$.
\end{proof}
\end{document}
